% Physical pages 3 and 4 of the concise study: the movement of the formulary,
% stated as an argument and not as a summary of what follows. The developed
% comparison of the three readings begins on the fifth page.

\runninghead{Themes and Movement}
\section{The Propers: Themes and Movement}
\zlabel{triptych:concise:themes:start}

This Mass sets a throne in heaven and a prisoner in a governor's hall, and it
sings of them as one King. Its chants give the slain Lamb power and divinity,
the King's Son the nations, and the Lord a seat for ever; its Gospel gives the
same King bound and questioned, saying that his kingdom is not of this world.
Between the two stands the sacrifice, which the Secret calls the victim of
mankind's reconciliation and the Preface the self-offering of an anointed
Priest who will hand over to the Father a kingdom without end.

\subsection{The throne and the nations}

The Introit begins where the Apocalypse has just shown a sealed book that no
one in heaven or on earth was worthy to open. The elders have sung to the Lamb,
``thou wast slain and hast redeemed us to God, in thy blood'' (Apoc 5:9), and
the angels take up the acclamation the antiphon keeps: \latin{Dignus est
Agnus, qui occisus est, accipere virtutem, et divinitatem}. The Lamb is worthy
because he was slain, and what he is worthy to receive includes divinity
itself. The antiphon adds, from John's greeting to the churches, ``to him be
glory and empire for ever and ever'' (1:6). Its verse is a prayer of giving
too, the first line of Psalm~71: ``Give to the king thy judgment, O God, and to
the king's son thy justice.'' Of the eternal Son, what can be received?

The Collect turns from the throne to the earth. God willed to restore all
things in his beloved Son, \latin{universorum Rege}, the King of all; its verb,
\latin{instaurare}, is the Vulgate's word for the summing up of all things in
Christ (Eph 1:10). It then names the condition of the peoples. They are
\latin{familiae gentium}, families of the nations, and they are
\latin{peccati vulnere disgregatae}: their scattering is a wound, and sin
inflicted it. What the prayer asks for them is subjection, \latin{eius
suavissimo subdantur imperio}, but it names the quality of the rule first: it
is most gentle.

\subsection{Translated into a kingdom that reaches every shore}

The Epistle says how anyone comes into that rule. Paul, writing in bonds, gives
thanks to the Father ``who hath delivered us from the power of darkness and
hath translated us into the kingdom of the Son of his love'' (Col 1:13). The
kingdom is not conquered; its subjects are carried into it. And Paul's hymn
then says who the King is: the image of the invisible God, the firstborn of
every creature, in whom thrones and dominations were created, the head of the
body, the Church, the firstborn from the dead, through whom the Father
reconciles all things, ``making peace through the blood of his cross''
(1:15--20).

The two chants after the Epistle measure that kingdom. The Gradual takes two
verses of Psalm~71 and leaves out those between them: ``he shall rule from sea
to sea, and from the river unto the ends of the earth''; ``all kings of the
earth shall adore him: all nations shall serve him'' (71:8, 11). The Alleluia
sings the second half of Daniel's vision of one like a son of man: ``his power
is an everlasting power that shall not be taken away: and his kingdom that
shall not be destroyed'' (Dan 7:14). It keeps the power as possessed and
lasting and passes over the words that begin the verse, ``And he gave him
power.'' Blessed Ildefonso Schuster reads the two chants together as ``the
double chant after the Epistle,'' a threefold catholicity: of place, from sea
to sea; of subjects, all kings and all nations; and of ages, a power that shall
not be taken away (\work{The Sacramentary} V, p.~203).

\subsection{The King before the governor}

Then the Gospel moves from every shore to one room. Jesus has been brought
early in the morning before the Pasch into the praetorium, which his accusers
will not enter, so that Pilate goes in and out between them and the prisoner
(Jn 18:28--29). ``Art thou the king of the Jews?'' Jesus answers with a
question of his own; Pilate replies that the prisoner's own nation has
delivered him up. Then the words on which the feast turns: ``My kingdom is not
of this world. If my kingdom were of this world, my servants would certainly
strive that I should not be delivered to the Jews: but now my kingdom is not
from hence.'' The Latin makes the denial three times, \latin{de hoc mundo},
\latin{ex hoc mundo}, \latin{hinc}. Pressed again, Jesus answers, ``Thou sayest
that I am a king. For this was I born, and for this came I into the world;
that I should give testimony to the truth. Every one that is of the truth
heareth my voice'' (18:36--37). The pericope stops before ``What is truth?''

Set beside the chants, the denial sharpens every claim the Mass has made. The
kingdom that rules from sea to sea is declared by a bound man to come from
somewhere else. The kings who are to adore him, and the nations whose
subjection the Collect asked, are claimed by a King who says that if his
kingdom were of this world his servants would strive for him, and that it is
not from hence. Yet the Postcommunion will ask that the faithful
\latin{militare}, fight as soldiers, under his standards.

\subsection{The sacrifice and the end of the kingdom}

The Offertory answers from the second psalm. The king set on Sion recites the
decree, ``Thou art my son, this day have I begotten thee,'' and the Lord goes
on, ``Ask of me, and I will give thee the Gentiles for thy inheritance, and the
utmost parts of the earth for thy possession'' (Ps 2:7--8). The antiphon sings
the second verse only, but it is spoken to the Son of the first, and the psalm
continues by bidding the kings of the earth understand and ``serve ye the Lord
with fear'' (2:10--11). Sung as the gifts are brought to the altar, it sets
the Son's asking at the moment of the offering.

The Secret names that offering by its effect. It is \latin{hostiam \dots\
humanae reconciliationis}, the victim of mankind's reconciliation, and the
victim is the Son himself, \latin{quem sacrificiis praesentibus immolamus}.
What it asks of him for all nations is \latin{unitatis et pacis dona}, the
gifts of unity and peace: the remedy of the Collect's wound. The proper
Preface then gives the argument of the whole feast in one sentence. The Father
anointed his only-begotten Son with the oil of gladness (Ps 44:8; Heb 1:9),
eternal Priest and King of all things, so that, offering himself on the altar
of the cross as a spotless and peace-making victim, he might accomplish the
mysteries of man's redemption, and, all creatures being subjected to his rule,
deliver to the Father's majesty ``an eternal and universal kingdom, a kingdom
of truth and life; a kingdom of holiness and grace, a kingdom of justice, love
and peace,'' in the English of the Lasance missal. Its verb,
\latin{traderet}, is Paul's ``when he shall have delivered up the kingdom to
God and the Father'' (1~Cor 15:24).

That handing over meets the chants that deny any end. The Alleluia has already
said that his power shall not be taken away, and the Communion now sings two
half-verses of Psalm~28: ``the Lord shall sit king for ever \dots\ the Lord
will bless his people with peace'' (28:10--11). Before them the psalm speaks
of a flood that the Lord makes to dwell, and between them of the strength he
gives his people; the antiphon leaves both unsung. The
Postcommunion, after \latin{immortalitatis alimoniam}, the food of
immortality, asks that those who fight under the King's standards may reign
with him for ever \latin{in caelesti sede}.

The formulary's movement thus raises three questions, and each of three
readings takes one as its own. The first begins from the Introit's verb and the Offertory's command: how the Lamb
receives what the Son always had. St~Hilary of Poitiers answers that the risen
Son asks back the glory he had before; St~Thomas Aquinas, that the Son holds
the nations by a right of birth and by a right won in prayer and passion;
St~Bede, that the Lamb who dies and the Lion who conquers are one Lord. The
second begins from the Gospel: where a kingdom not from hence is, and how it
claims kings and nations. St~Augustine answers that it is not of this world
but is in it; St~John Chrysostom, that its empire is from above and needs no
servants to defend it; St~Robert Bellarmine, that the kings are bidden to
subject themselves freely and will all bend at the judgment. The third begins
from the Preface: how a kingdom handed over can have no end. Augustine answers
that the kingdom handed over is the faithful, brought to the contemplation of
God; Hilary, that the Son delivers to God not his power but us, made his
kingdom by the glorifying of his body.
\zlabel{triptych:concise:themes:end}
